\chapter{Gestión del proyecto y seguimiento SCRUM}

Este capítulo recoge la gestión del proyecto con un enfoque SCRUM adaptado a un TFM individual. Se busca dejar trazabilidad clara de qué se planificó, qué se ejecutó, qué bloqueos aparecieron y qué decisiones técnicas se tomaron en cada iteración.

\section{Product Backlog}

\subsection{Backlog inicial}

El backlog de alto nivel se definió con los siguientes bloques:

\begin{enumerate}
  \item Prototipo web funcional con mapa, selección origen-destino y cálculo de rutas.
  \item Integración robusta de OSRM por perfiles (coche, bicicleta y a pie).
  \item Integración de transporte público con OTP y datos GTFS de Toledo.
  \item Desarrollo de pipeline reproducible de IA con el dataset LPMC.
  \item Entrenamiento y validación de un modelo base (XGBoost).
  \item Integración futura de inferencia en backend mediante endpoint específico.
  \item (WIP) Línea paralela de experimentación con DNN (Keras/TensorFlow en CPU).
  \item Consolidación de documentación técnica y memoria final.
\end{enumerate}

\subsection{Priorización de items}

La priorización siguió este criterio:

\begin{itemize}
  \item \textbf{Prioridad alta}: funcionalidades necesarias para disponer de un simulador operativo extremo a extremo.
  \item \textbf{Prioridad media}: mejoras de experiencia de usuario y visualización.
  \item \textbf{Prioridad evolutiva}: integración IA en producción y comparación ampliada de modelos.
\end{itemize}

\section{Histórico de sprints ejecutados}

\subsection{Sprint 1: Prototipo OSRM remoto}

\textbf{Objetivo}: validar el flujo mapa $\rightarrow$ backend $\rightarrow$ rutas por modo.

\textbf{Trabajo realizado}:
\begin{itemize}
  \item Arranque del frontend con React + Leaflet.
  \item Backend FastAPI mínimo para cálculo de rutas.
  \item Integración inicial con demoserver oficial de OSRM.
\end{itemize}

\textbf{Problema detectado}: el demoserver devuelve siempre la ruta en coche, lo que no permitía comparación entre \texttt{driving}, \texttt{cycling} y \texttt{foot}.

\textbf{Entregables}: base funcional de la aplicación y validación temprana de arquitectura cliente-servidor.

\subsection{Sprint 2: OSRM local multi-perfil}

\textbf{Objetivo}: obtener rutas diferenciadas y control total de infraestructura.

\textbf{Trabajo realizado}:
\begin{itemize}
  \item Paso intermedio por instancias públicas de \texttt{routing.openstreetmap.de}.
  \item Despliegue final de tres instancias OSRM locales en Docker.
  \item Configuración de puertos: 5000 (coche), 5001 (bici), 5002 (a pie).
  \item Consolidación del endpoint \texttt{/api/osrm/routes}.
\end{itemize}

\textbf{Entregables}: comparación modal estable y reproducible en entorno local.

\subsection{Sprint 3: GTFS Toledo}

\textbf{Objetivo}: incorporar red de transporte público urbano en la aplicación.

\textbf{Trabajo realizado}:
\begin{itemize}
  \item Carga y modelado del feed \texttt{GTFS\_Urbano\_Toledo}.
  \item Endpoints para paradas, rutas y horarios.
  \item Visualización de paradas y líneas en el frontend.
\end{itemize}

\textbf{Entregables}: capa GTFS operativa para exploración de red y horarios.

\subsection{Sprint 4: OTP multimodal}

\textbf{Objetivo}: habilitar rutas puerta a puerta con transporte público.

\textbf{Trabajo realizado}:
\begin{itemize}
  \item Construcción del \texttt{graph.obj} con OSM + GTFS.
  \item Endpoint \texttt{/api/otp/routes} con ordenación y paginación de itinerarios.
  \item Segmentación de legs (WALK/BUS) y exposición de metadatos de línea.
\end{itemize}

\textbf{Entregables}: integración completa de modo tránsito con desglose por segmentos.

\subsection{Sprint 5: Mejora visual y UX}

\textbf{Objetivo}: hacer más clara la comparación entre modos.

\textbf{Trabajo realizado}:
\begin{itemize}
  \item Estilos diferenciados por modo (colores y trazos).
  \item Visualización de segmentos OTP solo en modo tránsito.
  \item Mejoras de panel de información y selección de líneas.
\end{itemize}

\textbf{Entregables}: interfaz más legible y usable para análisis de escenarios.

\subsection{Sprint 6: LPMC exploración y preprocesado}

\textbf{Objetivo}: investigar el dataset LPMC y los trabajos de referencia (papers adjuntos) para preparar una base sólida de modelado.

\textbf{Trabajo realizado}:
\begin{itemize}
  \item Revisión técnica del dataset LPMC y de los papers de referencia incluidos en la carpeta \texttt{papers}.
  \item Scripts de exploración y preprocesado del dataset LPMC.
  \item Codificación de variables categóricas y limpieza de columnas.
  \item Split temporal train/test por \texttt{survey\_year}.
\end{itemize}

\textbf{Entregables}: datasets preprocesados reproducibles para entrenamiento de modelos.

\subsection{Sprint 7: XGBoost baseline}

\textbf{Objetivo}: entrenar un primer modelo útil y evaluable.

\textbf{Trabajo realizado}:
\begin{itemize}
  \item Entrenamiento de XGBoost multiclase.
  \item Escalado de variables numéricas y persistencia de artefactos.
  \item Validación por métricas de test e inspección de inferencias.
\end{itemize}

\textbf{Entregables}: \texttt{xgb\_lpmc\_baseline.joblib}, \texttt{xgb\_lpmc\_scaler.joblib} y resultados base para comparación.

\subsection{Sprint 8: Arranque de memoria y trazabilidad SCRUM}

\textbf{Objetivo}: iniciar la redacción formal de la memoria y consolidar el registro estructurado del trabajo realizado.

\textbf{Trabajo realizado}:
\begin{itemize}
  \item Inicio de la escritura de la memoria en LaTeX con estructura completa por capítulos.
  \item Propuesta y consolidación del índice actual del TFM.
  \item Registro del trabajo acumulado en formato SCRUM para seguimiento con tutor.
  \item Redacción inicial del capítulo de estado del arte.
\end{itemize}

\textbf{Entregables}: base de memoria preparada para redacción iterativa y seguimiento SCRUM trazable.

\subsection{Sprint 9: Integración de inferencia LPMC en simulador}

\textbf{Objetivo}: conectar el modelo de elección modal con backend y frontend para inferencia operativa por escenario.

\textbf{Trabajo realizado}:
\begin{itemize}
  \item Implementación de endpoint \texttt{/api/lpmc/predict} en FastAPI.
  \item Implementación de endpoint \texttt{/api/lpmc/debug-features} para inspección del vector de entrada al modelo.
  \item Integración en frontend de formulario de perfil usuario y visualización de probabilidades por modo.
  \item Ajuste de la estrategia de \texttt{household\_id}: reentrenamiento de una variante de modelo sin esta variable como feature (\texttt{nohh}), usándola solo para agrupación de folds.
  \item Conservación de artefactos legacy (con \texttt{household\_id}) como respaldo y conmutación por variable de entorno (\texttt{LPMC\_MODEL\_VARIANT}).
  \item Correcciones de interfaz (castellano, orden de bloques y legibilidad de resultados).
\end{itemize}

\textbf{Entregables}: flujo completo de inferencia (entrada de variables, cálculo de features de ruta, predicción y depuración de features).

\section{Sprint reviews y retrospectivas}

\subsection{Aprendizajes por iteración}

\begin{itemize}
  \item Validar pronto las limitaciones de servicios externos evita retrabajo (caso demoserver OSRM).
  \item Mantener infraestructura local mejora reproducibilidad y control experimental.
  \item Separar integración funcional y capa de IA reduce complejidad y facilita depuración.
  \item Mantener trazabilidad escrita del proyecto es clave para retomar el trabajo con rapidez tras varias semanas de pausa.
  \item Exponer endpoints de depuración acelera la validación del pipeline de inferencia en integración.
\end{itemize}

\subsection{Acciones de mejora}

\begin{itemize}
  \item Integrar inferencia en backend sin bloquear el servidor.
  \item Sustituir parámetros temporales hardcodeados de OTP por fecha/hora seleccionadas por usuario.
  \item Revisar la coherencia entre oferta real de transporte público y espacio de decisión del modelo: en trayectos cortos donde OTP no devuelve tramo BUS (itinerario íntegramente WALK), no debe habilitarse la alternativa \texttt{pt} en inferencia.
  \item Añadir trazabilidad explícita entre evidencias (scripts, notebooks, endpoints) y sprints.
  \item Estandarizar formato de revisión al cierre de cada iteración.
\end{itemize}

\section{Seguimiento operativo}

\subsection{Registro en SCRUM\_LOG}

El seguimiento detallado se mantiene en el documento vivo \texttt{latex/SCRUM\_LOG.md}, donde se registran:

\begin{itemize}
  \item estado de cada sprint,
  \item bloqueos,
  \item evidencia técnica,
  \item decisiones tomadas y próximos pasos.
\end{itemize}

\subsection{Trazabilidad con el repositorio}

Cada sprint debe poder vincularse con evidencias concretas:

\begin{itemize}
  \item archivos modificados,
  \item scripts ejecutables,
  \item artefactos de modelo,
  \item endpoints funcionales,
  \item resultados observables en interfaz.
\end{itemize}

\section{Próximo sprint propuesto}

\subsection{Objetivo}

Actualizar OTP y datos de transporte para 2026, eliminando el parche temporal de fecha/hora fija y utilizando fecha/hora de viaje seleccionadas por el usuario.

\subsection{Alcance y criterio de cierre}

\begin{itemize}
  \item Feed GTFS actualizado (2026) y \texttt{graph.obj} regenerado.
  \item Selector de fecha y hora de viaje integrado en frontend.
  \item Paso de \texttt{date}/\texttt{time} desde UI hasta \texttt{\_build\_otp\_params} en backend.
  \item Comprobación de itinerarios OTP con fecha/hora variable en Toledo.
  \item Evidencia registrada en \texttt{SCRUM\_LOG.md} y en la memoria.
\end{itemize}

\subsection{Plan a dos sprints vista}

La línea de investigación DNN (Keras/TensorFlow en CPU, ajuste de depth/width y dropout) se planifica para dos sprints posteriores al de actualización OTP, una vez estabilizado el flujo de rutas con fecha/hora seleccionables. En paralelo, se propone una tarea de robustez metodológica en inferencia modal: restringir el conjunto de alternativas a modos realmente disponibles en cada consulta. En particular, si el planificador no devuelve transporte público efectivo (sin tramo BUS), la elección debe realizarse solo entre \texttt{walk}, \texttt{cycle} y \texttt{drive}, evitando clasificar como \texttt{pt} desplazamientos que operativamente se resuelven andando.
